\begin{document}
\title{ACTS-Attack: Measuring Cryptographic Attack\\Construction, Admission, and Attribution}
\author{}
\maketitle

\begin{abstract}
Security-agent evaluations report attack success, but construction, evaluator
admission, and execution attribution require different evidence. ACTS-Attack
studies these requirements through cryptographic recovery programs,
target-effect controls, and adversarial evidence bundles. A matched study
contains 768 representation trials and 84 controls. On 128 instances at the
primary observation-complexity level (L3), the
base64--hexadecimal success contrast is $-0.8$ percentage points under the
declared contract, within the pre-specified $\pm10$ point margin for two
families and one route. Applying a contract-mismatched source filter to
unchanged programs moves it to $-32.8$ points; output-preserving unused-import
removal changes the filtered contrast to $+3.9$ points. Separate construction
studies distinguish 12/12 supplied operations from 8/12 accepted prospective
sessions, locate observation and implementation failures, and demonstrate
protected-payload release through a supplied AES-CTR chain. Constructed
attribution tests show that candidate correctness, record integrity, and
signatures alone can leave execution claims unsupported. A claim-specific
policy undergoes 11{,}256 bounded mutation evaluations without an observed
soundness violation; configured in-toto agreement tests its implementation.
Together, these studies provide program-level measurements and concrete
failure mechanisms for interpreting cryptographic attack claims. Their
cohorts retain separate denominators: construction spans distinct local
interfaces, admission sensitivity remains local to H1-R3, and attribution
depends on the stated trust model. Retaining decision stages and execution
bindings makes these distinct claims auditable.
\end{abstract}
\begin{IEEEkeywords}
AI agents, cryptographic attacks, exploit construction, evaluation validity,
verification, provenance
\end{IEEEkeywords}
